% ── §R1: Omicron Spike Protein — Primary Analysis ──
\section{Detailed Results: SARS-CoV-2 Omicron Spike Analysis}
\label{sec:results-spike}

Before presenting cross-protein transfer results on the PSICOV150 benchmark,
we present the complete analysis of our primary target: the SARS-CoV-2
Omicron variant Spike protein.  This serves both as a biological case study
and as a demonstration of the full pipeline applied to a single, deeply
characterised protein family.

\subsection{Dataset Characteristics}
\begin{table}[ht]
	\centering
	\caption{Omicron Spike dataset properties.}
	\label{tab:spike-data}
	\begin{tabular}{lr}
		\toprule
		Property                                 & Value                                          \\
		\midrule
		Sequences                                & 1{,}299                                        \\
		Alignment width                          & 1{,}276 columns                                \\
		Unique sequences                         & 299 (587 identical copies of dominant variant) \\
		Variable positions ($H > 0.3$)           & 21 of 1{,}276                                  \\
		Co-evolving pairs (mut.-only MI $> 0.1$) & 10                                             \\
		Forbidden rules (flipped QM)             & 490                                            \\
		Conserved positions                      & 1{,}255                                        \\
		Mean mutations per sequence              & 11.1                                           \\
		\bottomrule
	\end{tabular}
\end{table}

The extreme conservation of the Spike protein (98.4\% of positions invariant)
reflects the structural constraints imposed by its trimeric assembly and
receptor-binding function.  Only 21 positions carry appreciable evolutionary
variability, concentrated in the N-terminal domain (NTD) and receptor-binding
domain (RBD) that are targets of antibody-mediated immune pressure.

\subsection{The Ten Co-Evolving Position Pairs}
\begin{table}[ht]
	\centering
	\caption{All ten co-evolving position pairs ranked by mutation-only mutual information.}
	\label{tab:spike-pairs}
	\begin{tabular}{clcccc}
		\toprule
		Rank & Pair       & Full MI & Mut.-only MI   & Sep. & Reference \\
		\midrule
		1    & (495, 498) & 0.039   & \textbf{0.871} & 3    & R--G      \\
		2    & (448, 454) & 0.027   & 0.834          & 6    & G--L      \\
		3    & (488, 498) & 0.032   & 0.822          & 10   & F--G      \\
		4    & (442, 454) & 0.000   & 0.811          & 12   & K--L      \\
		5    & (442, 448) & 0.016   & 0.728          & 6    & K--G      \\
		6    & (212, 215) & 0.398   & 0.398          & 3    & V--G      \\
		7    & (215, 216) & 0.757   & 0.377          & 1    & G--R      \\
		8    & (212, 216) & 0.377   & 0.377          & 4    & V--R      \\
		9    & (210, 215) & 0.753   & 0.238          & 5    & N--G      \\
		10   & (210, 212) & 0.177   & 0.177          & 2    & N--V      \\
		\bottomrule
	\end{tabular}
\end{table}

A critical observation is that the top five pairs by mutation-only MI (the
442--498 cluster) have very low full MI (0.000--0.039), indicating that their
covariation is driven by mutations away from the reference rather than by the
reference itself.  This is the signature of lineage-specific covariation:
different viral variants carry different combinations at these positions,
producing strong mutation correlation without necessarily implying physical
contact.

\subsection{Prime Implicants and Essential Rules}
Quine--McCluskey minimisation over ten co-evolving pairs produces 36 distinct
prime implicants, of which exactly 2 are essential (the minimal irredundant
core):
\begin{align}
	 & \textbf{Rule E1:} \quad \text{pos}212 = \text{G} \;\wedge\; \text{pos}216 = \text{R}, \\
	 & \textbf{Rule E2:} \quad \text{pos}210 = \text{K} \;\wedge\; \text{pos}215 = \text{G}.
\end{align}
Each essential rule fires in only one of 1{,}299 sequences, reflecting the
rarity of the specific double-mutant combination.  The remaining 34 non-essential
implicants cover the broader coevolutionary landscape at each pair.

\subsection{Flipped Circuit: Forbidden Combinations}
Applying Quine--McCluskey to the flipped maps (ON = never observed, OFF =
observed, DC = reference) produces 490 forbidden rules across 117 unique
residue-pair identities.  These represent negative-selection constraints:
combinations that are never sampled across 1{,}299 sequences despite adequate
sampling depth at both positions.

Three-dimensional validation reveals that forbidden pairs at the NTD
positions (210, 215) and (212, 215) sit at contact-compatible distances
(enrichment $2.57\times$ and $2.71\times$ respectively), consistent with
steric incompatibility driving their absence.  Forbidden pairs in the
442--498 cluster show no such enrichment, supporting their interpretation as
lineage-covariation artefacts rather than structural constraints.

\subsection{Structural Validation Against 299 Homology Models}
We mapped each alignment column to the corresponding residue in the best-GMQE
SWISS-MODEL structure (identity mapping verified for all 299 unique
sequences).  Among the ten co-evolving pairs:

\begin{table}[ht]
	\centering
	\caption{3D validation of co-evolving pairs against SWISS-MODEL structures.}
	\label{tab:spike-3d}
	\begin{tabular}{lcccl}
		\toprule
		Pair       & Contact frac. & Enrichment   & Median \AA{} & Region         \\
		\midrule
		(407, 410) & 88.6\%        & 6.97$\times$ & 6.7          & RBD            \\
		(212, 215) & 35.4\%        & 2.78$\times$ & 8.5          & NTD            \\
		(210, 215) & 29.3\%        & 2.30$\times$ & 11.0         & NTD            \\
		(500, 503) & ---           & 11.3$\times$ & 5.3          & RBD            \\
		(503, 507) & ---           & 11.3$\times$ & 5.0          & RBD            \\
		(454, 495) & ---           & 5.97$\times$ & 8.0          & RBD (DCA-only) \\
		\bottomrule
	\end{tabular}
\end{table}

Seven pairs achieve enrichment $\geq 2\times$ with $p < 10^{-6}$, confirming
that the Karnaugh-map framework detects genuine structural signal.  However,
the strongest full-MI pair (373, 378; MI $= 0.807$) shows zero enrichment,
demonstrating that high covariation does not necessarily imply physical
proximity when the underlying cause is shared phylogenetic history rather
than structural constraint.

\subsection{Perplexity Ratio as a Structural Predictor}
Among all sequence-derived metrics evaluated on the Spike, the perplexity
ratio achieves the highest Spearman correlation with native-contact fraction:
\begin{equation}
	\rho(\text{ratio}, \text{contact}) = 0.670, \quad p = 0.006.
\end{equation}
Near-deterministic pairs (ratio $> 1.5$) are enriched for contacts at
$1.80\times$, while low-ratio pairs show no enrichment.  This suggests that
perplexity ratio captures a property of coevolution---near-determinism of the
amino-acid relationship---that raw MI does not distinguish from noisy
covariation.

\subsection{Conserved Core Analysis}
Of the 1{,}276 alignment columns, 1{,}255 are conserved ($H < 0.3$).  These
columns encode the structural core of the Spike protein, including the fusion
peptide, heptad repeats, and transmembrane domain.  While they generate no
coevolutionary signal (their entropy is too low for meaningful MI), they
carry structural information through their absolute constraint: any mutation
at these positions is likely deleterious and rapidly purged.  Our complete
circuit incorporates this information through the conservation-type feature
dimension, which enriches cons/cons column pairs for native contacts by
$3.2\times$.
